\emph{Draft, 2026-07-10. Qwen3-1.7B/8B. All numbers reproducible from
\texttt{scripts/} and \texttt{data/}; see \href{reproduce.md}{"How to
reproduce"} and the \href{findings.md}{working findings log}. An
\href{explorer.md}{interactive explorer} animates the three central
results (declension, injection, syncretism) and, for readers new to
grammatical case, explains the linguistic analogy.}

\subsection{Abstract}\label{abstract}

We study how an open-weights transformer (Qwen3) represents a
\textbf{relational operation} applied to an entity ---
\emph{currency-of}, \emph{capital-of}, \emph{language-of},
\emph{demonym-of}, \emph{continent-of} applied to a country. We find
that the operation is carried by a \textbf{manipulable operator
direction} in the residual stream that is largely \textbf{separable from
the operand}. Adding \texttt{v(op\_B)\ -\ v(op\_A)} to a prompt
performing \texttt{op\_A} converts the model\textquotesingle s answer to
\texttt{op\_B} for all 20 ordered operator pairs (matched-norm random
controls null), and the direction \textbf{generalizes to operands never
used to build it}. A two-way analysis of variance factorizes the
representation as
\texttt{H{[}operand,\ operator{]}\ =\ \ensuremath{\mu}\ +\ operand\ +\ operator\ +\ interaction},
\textasciitilde90\% additive with a \textasciitilde10\% \textbf{fusion}
interaction, operator and operand occupying largely orthogonal
subspaces; the operation-dominant representation \emph{emerges along the
sequence} (operand-dominant at the entity token \ensuremath{\to} operator-dominant at
the query token). Strikingly, the operation is \textbf{separable from
its surface realization}: \emph{language-of} and \emph{demonym-of} are
distinct operator directions even though both emit "Italian". We frame
this with a morphological metaphor --- the model \textbf{declines} a
concept by its functional role (case), with \textbf{syncretism} (shared
surface form) and \textbf{fusion} (non-additive stem+ending) as in a
Latin declension. Testing generalization across domains, the
factorization is \textbf{specific to relational retrieval}: for
\textbf{arithmetic} (+, \ensuremath{\times}, -) and \textbf{comparison-logic} operators,
the interaction term is 2--4\ensuremath{\times} larger and the operator \textbf{fails to
generalize to held-out operands}, consistent with arithmetic being
computed by a "bag of heuristics" rather than a linear operator. A
control that instead varies a \emph{linear numeric parameter} (add-N)
rather than the function generalizes better and is near-collinear,
reconciling a contrary recent report (Christ et al. 2025). Finally, the
structure is \textbf{causal, not a readout advantage}: the Jacobian-lens
("J-space") readout is not more legible than the raw residual on any of
our metrics under matched controls.

\subsection{1. Introduction}\label{1-introduction}

A line of work shows that a \emph{task} or \emph{relation} can be
captured by a single addable vector in an LLM\textquotesingle s residual
stream (task vectors, Hendel et al. 2023; function vectors, Todd et al.
2024) and that many relations are approximately linear operators (LRE,
Hernandez et al. 2024; Merullo et al. 2024). What has not been
characterized is the \textbf{joint structure of operator and operand}:
whether the operation \emph{factorizes} from its argument, whether that
factorization \emph{generalizes}, whether the operation is separable
from the \emph{word it produces}, and whether any of this holds
\emph{beyond} relational facts, in arithmetic and logic. We answer these
on Qwen3, and give a morphological reading --- relational computation as
\textbf{declension} --- that predicts the specific asymmetries we
observe.

\subsection{2. Related work and what is new
here}\label{2-related-work-and-what-is-new-here}

\textbf{Operation as an addable vector.} Task Vectors (Hendel et al.,
EMNLP 2023) and Function Vectors (Todd et al., ICLR 2024) extract one
vector per ICL task and add/patch it to trigger the task. In-Context
Vectors (Liu et al., ICML 2024) and ActAdd (Turner et al. 2023) steer a
task/style similarly. These bundle operator \emph{and} operand into one
unfactored "task" direction and (for FVs) show \emph{partial} cross-task
arithmetic without null controls. We contribute the
\textbf{factorization} (operator \ensuremath{\oplus} operand), an \textbf{exhaustive
all-pairs swap with matched-norm nulls}, and a \textbf{held-out-operand
generalization} test.

\textbf{Relation as a linear operator.} LRE (Hernandez et al., ICLR
2024) fits a full affine map \texttt{W\textperiodcentered{}s\ +\ b} per relation. That
represents the relation as a \emph{monolithic} operator; we show a
\textbf{low-dimensional additive operator \emph{direction}} that is
arithmetically composable and \textbf{factorable from the operand}, and
we quantify the operand/operator/interaction variance split, which LRE
does not.

\textbf{Additive factual recall.} Summing Up the Facts (Chughtai et al.
2024) decomposes factual recall into additive \emph{circuits} (subject
vs relation contributions). Ours is a \textbf{representational two-way
ANOVA of the residual state} with an explicit, measured
\textbf{interaction/fusion term} --- a different object.

\textbf{Attribute subspaces / reading position.} RAVEL (Huang et al.,
ACL 2024) disentangles multiple \emph{attributes of one entity} into
subspaces (operators sharing an operand); we measure the orthogonal
axis, \textbf{operator vs operand}. The operand\ensuremath{\to}operator shift along the
sequence restates the subject-enrichment \ensuremath{\to} attribute-extraction dynamics
of Geva et al. (EMNLP 2023) in operator/operand terms --- our least
novel point, included as confirmation.

\textbf{Arithmetic and logic.} Arithmetic in LLMs is reported as a "bag
of heuristics" (Nikankin et al. 2025) and Fourier/helical procedures
(Nanda et al. 2023; Zhong et al. 2023; Kantamneni \& Tegmark 2025), with
the operand \emph{value} cleanly linear (Gurnee \& Tegmark 2024). Truth
is a steerable direction (Marks \& Tegmark 2023) and logical-operator
heads exist (Hong et al. 2024). No prior work reports a
\textbf{cross-domain operator/operand factorization}; we provide one and
show it \textbf{breaks down} for arithmetic and logic --- itself a
result.

\textbf{Genuinely new here:} (1) operation-vs-realization dissociation
(language \ensuremath{\neq} demonym as directions despite identical output); (2) the
operator\ensuremath{\times}operand ANOVA with a named fusion term; (3) the
declension/case-paradigm framing; (4) all-pairs swap with matched-norm
nulls; (5) operator/operand subspace orthogonality; (6) the cross-domain
generalization test that separates relational (clean) from
arithmetic/logic (entangled).

\subsection{3. Method}\label{3-method}

\textbf{Setup.} Qwen3-1.7B and 8B, run on a single AMD Strix Halo APU
(see \href{setup.md}{setup}; an int8 path keeps larger models in range).
A relation is rendered into a uniform template
\texttt{"The\ \{op\}\ of\ \{a\}\ is"}. Reading position is the query
token unless stated; the depth "workspace" band (38--92\% of layers) is
where the J-space paper locates the global workspace and where we build
operator directions.

\textbf{Operator direction.} For operator \texttt{k},
\texttt{v(k)\ =\ mean\_operand{[}\ h(operand,\ k)\ -\ mean\_op\ h(operand,\ \textperiodcentered{})\ {]}}
at the workspace layers --- the deviation of an operator from the
operand\textquotesingle s average over operators, averaged over operands
(the function-vector construction;
\texttt{scripts/op\_core.py:op\_dirs}).

\textbf{Swap and controls.} Efficacy of \texttt{k\_A\ \ensuremath{\to}\ k\_B} is the
change in \texttt{logit(answer\_B)\ -\ logit(answer\_A)} at the query
position when \texttt{\ensuremath{\alpha}\textperiodcentered{}(v(k\_B)\ -\ v(k\_A))} is added over the band,
versus a \textbf{matched-norm random} direction. Answers are scored by
their distinguishing token (bare digit for arithmetic, since Qwen3
splits " 3" into {[}space, 3{]}); a tokenization guard drops operator
pairs that collide at that token.

\textbf{Factorization.}
\texttt{H{[}operand,\ operator{]}\ =\ \ensuremath{\mu}\ +\ operand(o)\ +\ operator(k)\ +\ interaction};
we report the variance share of each term and the principal angles
between the operand- and operator-subspaces
(\texttt{scripts/op\_core.py:factorize}).

\textbf{Generalization.} Build \texttt{v(op)} from half the operands;
test the swap on the held-out half. This distinguishes a genuine
operator from interpolation among the examples used to build it.

\textbf{J-space controls.} We repeat readout-geometry in the J-lens
readout \texttt{unembed(J\textperiodcentered{}h)} vs the logit-lens readout
\texttt{unembed(h)}, and re-run efficacy with \textbf{spectrum-matched
random-projection} and \textbf{permuted-vocabulary} null lenses
(\texttt{scripts/control\_lens.py}).

\subsection{4. Results}\label{4-results}

\subsubsection{4.1 Relational operators are manipulable
directions}\label{41-relational-operators-are-manipulable-directions}

All 20 ordered operator swaps flip the answer (mean swap \textbf{+21}
logit units on 1.7B, \textbf{+25} on 8B) while the matched-norm random
control is \textasciitilde0. Representative (1.7B):

{
\begin{longtable}[]{@{}lrrr@{}}
\toprule\noalign{}
swap & clean (from wins) & +operator & +random \\
\midrule\noalign{}
\endhead
\bottomrule\noalign{}
\endlastfoot
currency \ensuremath{\to} capital & -7.9 & \textbf{+31.5} & -1.2 \\
capital \ensuremath{\to} currency & -12.0 & \textbf{+33.0} & -2.5 \\
capital \ensuremath{\to} language & -8.5 & \textbf{+32.5} & +3.0 \\
continent \ensuremath{\to} language & -1.9 & \textbf{+23.0} & +1.4 \\
\end{longtable}
}

\pandocbounded{\includegraphics[width=\linewidth,keepaspectratio,alt={All-pairs operator swap: clean (baseline, "from" wins, blue) vs. injecting v(to)-v(from) ("to" wins, red). 20/20 pairs flip sign; matched-norm random control \ensuremath{\approx} 0.}]{op_swap}}

\emph{Injecting the operator difference \texttt{v(to)\ -\ v(from)} flips
every one of the 20 ordered pairs: clean (left, "from" wins, blue) \ensuremath{\to}
swapped (right, "to" wins, red). The matched-norm random control moves
nothing.}

\subsubsection{4.2 The representation factorizes into operator \ensuremath{\oplus}
operand}\label{42-the-representation-factorizes-into-operator--operand}

Two-way ANOVA at a mid-workspace layer (1.7B / 8B):

{
\begin{longtable}[]{@{}lrrr@{}}
\toprule\noalign{}
read position & operand & operator & interaction (fusion) \\
\midrule\noalign{}
\endhead
\bottomrule\noalign{}
\endlastfoot
query token & 5\% / 6\% & \textbf{86\% / 82\%} & 9\% / 13\% \\
entity token & \textbf{59\% / 55\%} & 31\% / 34\% & 9\% / 11\% \\
\end{longtable}
}

\textasciitilde90\% additive; operand/operator subspaces largely
orthogonal (principal angles 41--82\textdegree{}). The operation- dominant
representation \textbf{emerges along the sequence}: the entity enters
operand-dominant and is \emph{declined} into an operator-marked form by
the query position (the reading-position control rules out template
echo; the causal swaps are the stronger control).

\pandocbounded{\includegraphics[width=\linewidth,keepaspectratio,alt={PCA of the 60 workspace vectors H{[}operand, operator{]}, colored by operator. At the country token the cases intermix (operand-organized, stem 59\%); at the query token they separate into clean case clusters (operator-organized, case 86\%). The faint web links each country\textquotesingle s five case-forms --- short at the country token, splayed at the query token. Bottom: the variance split, 1.7B and 8B.}]{op_geometry}}

\emph{Where operand and operator live. Colored by operator throughout:
at the \textbf{country token} the colors intermix (the cloud is
organized by operand); at the \textbf{query token} the same points
separate into five case clusters. Each country\textquotesingle s five
case-forms (faint web) start together and are pulled apart --- the
concept is \textbf{declined} along the sequence.}

\subsubsection{4.3 Operation \ensuremath{\neq} realization
(syncretism)}\label{43-operation--realization-syncretism}

\emph{language-of} and \emph{demonym-of} emit the \textbf{same word}
(Italian) yet are \textbf{distinct operator directions}; the swap
between them is weak precisely because they share an exponent. A
\textbf{pure desinence} built from the shared-output pairs
(\texttt{mean{[}h(language)\ -\ h(demonym){]}}, exponent cancelled)
still installs the relation causally (1.7B: -3.7 \ensuremath{\to} +12.5). The
syncretism appears at the \emph{exponent} level, not the \emph{case}
level --- as in Latin declension.

\pandocbounded{\includegraphics[width=\linewidth,keepaspectratio,alt={Left: operator-direction cosines --- all five operators are distinct directions (no off-diagonal near \ensuremath{\pm}1); language and demonym, boxed, both emit "Italian" yet differ (cos -0.26). Right: the pure desinence, built where the two share an output word, still installs the relation causally (-3.6 \ensuremath{\to} +12.5).}]{op_syncretism}}

\emph{Operation \ensuremath{\neq} realization. Every operator is a distinct direction
--- including language and demonym (boxed), which emit the identical
word "Italian". A desinence built precisely where the two \textbf{share}
their output word (so the word cancels) still installs the relation: the
case is separable from the word that realizes it.}

\subsubsection{4.4 Operators generalize to held-out
operands}\label{44-operators-generalize-to-held-out-operands}

\texttt{v(op)} built on 6 operands and applied to the 6 held-out
operands still flips \textbf{20/20} swaps (mean \textbf{+19} on 1.7B)
--- a genuine, transferable operator, not interpolation.

\subsubsection{4.5 The factorization is domain-specific (arithmetic and
logic do
not)}\label{45-the-factorization-is-domain-specific-arithmetic-and-logic-do-not}

Extending the identical pipeline to arithmetic (+, \ensuremath{\times}, -) and
comparison-logic operators, on 1.7B:

{
\begin{longtable}[]{@{}lrrll@{}}
\toprule\noalign{}
domain & operator variance & interaction (fusion) & held-out
generalization & swap vs random \\
\midrule\noalign{}
\endhead
\bottomrule\noalign{}
\endlastfoot
\textbf{relational} & 86\% & 9\% & \textbf{20/20} flip & +21 vs
\textasciitilde0 \\
\textbf{arithmetic} & 55\% & 23\% & 2/6 (\ensuremath{\leq} random) & +0.4 vs +0.3 \\
\textbf{logic (compare)} & 33\% & 34\% & 2/6 (\ensuremath{\leq} random) & +1.2 vs
+0.2 \\
\end{longtable}
}

The gradient replicates and sharpens at \textbf{8B}: relational held-out
generalization 20/20 (interaction 13\%); arithmetic interaction 25\%,
held-out 1/6; logic interaction 45\%, held-out \textbf{0/6} (swap below
random). Monotone at both scales: relational \ensuremath{\gg} arithmetic \textgreater{}
logic.

Relational operators factorize cleanly and generalize; arithmetic and
logical operators are \textbf{entangled with their operands} (2--4\ensuremath{\times} the
interaction) and do \textbf{not} form a transferable direction ---
consistent with arithmetic being a bag of heuristics / Fourier
computation rather than a linear operator. The \textbf{held-out
generalization test is the discriminator} between a real operator and
memorized interpolation.

\subsubsection{4.6 Reconciling with a contrary report: add-N vs.
two-operand
arithmetic}\label{46-reconciling-with-a-contrary-report-add-n-vs-two-operand-arithmetic}

Christ et al. (2025) report the \emph{opposite} sign for arithmetic: an
operator built from an \textbf{add-N} relation (a fixed addend
\texttt{N} applied to a single number) \textbf{does} generalize to
held-out relations. The two results do not conflict --- they cut
arithmetic along different axes. In add-N the operator \textbf{is the
addend \texttt{N}} over one numeric operand, so "generalizing across
\texttt{N}" is interpolation along the number line (the operand
\emph{value} is itself linear; Gurnee \& Tegmark 2024), not transfer of
a \emph{function}. Our \texttt{+\ \ensuremath{\times}\ -} cut varies the \textbf{function}
over two operands. Running both cuts on Qwen3
(\texttt{op\_core.py\ -\/-domain\ arith\_addN} vs \texttt{arithmetic};
\texttt{scripts/operator\_collinearity.py}) places add-N exactly between
relations and the genuine-function cut:

{
\begin{longtable}[]{@{}llr@{}}
\toprule\noalign{}
cut & held-out generalization (1.7B / 8B) & operator-set collinearity
(top-1 var) \\
\midrule\noalign{}
\endhead
\bottomrule\noalign{}
\endlastfoot
relations (5 operations) & 20/20 \textperiodcentered{} 20/20 & 0.39 (spread paradigm) \\
\textbf{add-N} (operator = addend) & 5/12 \textperiodcentered{} 6/12 (+0.10 / +0.05) &
\textbf{0.76 (most 1-D / number-line)} \\
\texttt{+\ \ensuremath{\times}\ -} (operator = function) & 2/6 \textperiodcentered{} 1/6 (-0.43 / -0.55) &
0.64 \\
\end{longtable}
}

add-N generalizes \textbf{better} than \texttt{+\ \ensuremath{\times}\ -} (positive vs.
clearly negative) and its operator directions are the \textbf{most
collinear} of the three families --- 76\% of the operator-set variance
on a single line --- consistent with add-N being a linear numeric family
rather than a set of distinct operations. This reproduces Christ et
al.\textquotesingle s positive (their cut varies a linear parameter)
alongside our negative (ours varies the function) with no contradiction.
The caveat is quantitative: Qwen3 BPE restricts us to single-digit
results, so the add-N grid is small (5 operands) and its generalization
is weak in absolute terms --- the \textbf{ordering} (relations \ensuremath{\gg} add-N
\textgreater{} \texttt{+\ \ensuremath{\times}\ -}), not the magnitude, is the claim.

\subsubsection{4.7 This is causal structure, not a J-space readout
advantage}\label{47-this-is-causal-structure-not-a-j-space-readout-advantage}

On four readout comparisons (bridge-entity pass@k, a surface-form logit
difference, a number probe, and a concept-plane trajectory), the J-lens
readout shows \textbf{no advantage} over the logit-lens under matched
controls; a spectrum-matched random projection reads number as well as
the fitted J-lens, and the concept-plane "channel" does not survive a
random-plane baseline (see \href{findings.md}{findings \S{}Part 1}). The
operator/operand structure lives in the \textbf{causal} organization of
the residual, and is \textbf{not} more legible in the J-space than in
the raw stream (its 1.7B readout-geometry difference does not replicate
at 8B). We therefore make no claim of a special readable subspace; the
contribution is the operator/operand \emph{causal} factorization.

\subsection{5. Limitations}\label{5-limitations}

\begin{itemize}
\tightlist
\item
  \textbf{Scale.} 1.7B and 8B; no 32B lens yet. Relational linearity
  holds for a subset of relations even in-domain (cf. LRE
  \textasciitilde48\%); our five relations are hand-chosen and
  English-only.
\item
  \textbf{Arithmetic coverage.} Qwen3 BPE splits multi-digit numbers, so
  we restrict to single-digit results; the arithmetic null is on that
  regime and at 1.7B/8B, where arithmetic competence is itself limited.
\item
  \textbf{Reading-position result} restates known
  subject-enrichment\ensuremath{\to}attribute-extraction dynamics.
\item
  \textbf{No J-space claim.} The readout-geometry difference is
  1.7B-specific; the paper stands on the causal factorization, not on a
  J-lens readout advantage.
\end{itemize}

\subsection{6. Conclusion}\label{6-conclusion}

An LLM represents a relational operation as a low-dimensional,
causally-manipulable operator direction that factorizes from the
operand, generalizes to unseen operands, and separates the operation
from the word that realizes it --- a structure well described as
\textbf{declension}. This factorization is not universal: it is clean
for relational retrieval and breaks down for arithmetic and logic, whose
operators are entangled with their operands. The structure is causal; it
is not a more-readable subspace.

\subsection{References}\label{references}

Christ et al. 2025, \emph{{[}title TK{]}}, arXiv:2510.26543 \textperiodcentered{} Chughtai
et al. 2024, \emph{Summing Up the Facts}, arXiv:2402.07321 \textperiodcentered{} Geva et al.
2023, \emph{Dissecting Recall of Factual Associations}, EMNLP \textperiodcentered{} Gurnee
\& Tegmark 2024, \emph{Language Models Represent Space and Time}, ICLR \textperiodcentered{}
Hendel et al. 2023, \emph{In-Context Learning Creates Task Vectors},
EMNLP Findings, arXiv:2310.15916 \textperiodcentered{} Hernandez et al. 2024,
\emph{Linearity of Relation Decoding in Transformer LMs} (LRE), ICLR,
arXiv:2308.09124 \textperiodcentered{} Hong et al. 2024, \emph{A Implies B: Circuit Analysis
for Propositional Logic}, arXiv:2411.04105 \textperiodcentered{} Huang et al. 2024,
\emph{RAVEL}, ACL, arXiv:2402.17700 \textperiodcentered{} Kantamneni \& Tegmark 2025,
\emph{Language Models Use Trigonometry to Do Addition}, arXiv:2502.00873
\textperiodcentered{} Liu et al. 2024, \emph{In-Context Vectors}, ICML, arXiv:2311.06668 \textperiodcentered{}
Marks \& Tegmark 2023, \emph{The Geometry of Truth}, arXiv:2310.06824 \textperiodcentered{}
Merullo et al. 2024, \emph{Language Models Implement Simple
Word2Vec-style Vector Arithmetic}, NAACL, arXiv:2305.16130 \textperiodcentered{} Nanda et
al. 2023, \emph{Progress Measures for Grokking}, ICLR, arXiv:2301.05217
\textperiodcentered{} Nikankin et al. 2025, \emph{Arithmetic Without Algorithms}, ICLR,
arXiv:2410.21272 \textperiodcentered{} Todd et al. 2024, \emph{Function Vectors in Large
Language Models}, ICLR, arXiv:2310.15213 \textperiodcentered{} Turner et al. 2023,
\emph{Activation Addition (ActAdd)}, arXiv:2308.10248 \textperiodcentered{} Zhong et al.
2023, \emph{The Clock and the Pizza}, NeurIPS, arXiv:2306.17844.
